\documentclass[10pt,a4paper]{amsart}
\usepackage[margin=19mm]{geometry}
\usepackage{amsmath,amssymb,mathtools}
\usepackage[scr=rsfs,cal=euler]{mathalfa}
\usepackage{xfrac}
\usepackage[unicode,hidelinks,bookmarks=false]{hyperref}
\newcommand{\ie}{i.e.\ }
\newcommand{\cf}{cf.\ }
\newcommand{\resp}{respectively\ }
\newcommand{\Subsection}{\subsection}
\providecommand{\nobreakdash}{\nobreak\hbox{-}}
\setlength{\emergencystretch}{2em}
\multlinegap0pt
\usepackage{bm}
\usepackage{bbm}
\usepackage{dsfont}
\usepackage{xspace}
\usepackage{cancel}
\usepackage{mathtools}
\usepackage{amsthm}
\makeatletter
\@ifundefined{theorem}{\newtheorem{theorem}{Theorem}}{}
\@ifundefined{lemma}{\newtheorem{lemma}[theorem]{Lemma}}{}
\@ifundefined{proposition}{\newtheorem{proposition}[theorem]{Proposition}}{}
\makeatother
\title[ReLU$^k$ Neural de Rham Complexes]{ReLU$^k$ Neural de Rham Complexes}
\author{Kaibo Hu, Jindong Wang, Jinchao Xu}
\date{Source: arXiv:2607.22478v1 (CC BY 4.0)}
\begin{document}
\maketitle

\noindent\emph{Source and licence.} ReLU$^k$ Neural de Rham Complexes. Version arXiv:2607.22478v1 (CC BY 4.0), distributed under the Creative Commons Attribution 4.0 International licence, as recorded in the arXiv licence metadata for this exact version and shown on the abstract page https://arxiv.org/abs/2607.22478v1. Full rights notice: licenses/arxiv-2607.22478-CC-BY-4.0.txt.\par\smallskip

\noindent\textit{Editorial scope.} Counted unit: the Theorem whose source label is \texttt{thm:exact} in the article (source0.tex lines 528--539, source label \texttt{thm:exact}), delivered together with the article's own complete original proof of it (source0.tex lines 541--616, one \texttt{proof} environment of 76 lines). No mathematics is written, repaired or re-derived: statement and proof are the article's own text line for line. Required context retained uncounted: prop:inclusion (lines 204--216); lem:induction (lines 272--282); lem:koszul (lines 467--488). Every definition, display and label of those lines is retained unchanged. Omitted from this package and not redistributed: 1--203, 217--271, 283--466, 489--527, 540--540, 617--1027. Nothing inside a retained excerpt is omitted or altered: every retained body is delivered exactly as the article's lines are written, so no omission receipt of any kind accompanies this package. Citations: the retained excerpts carry no citation command at all, so no bibliography entry of the article is transcribed and no reference is redistributed. Reference adaptation: none was needed and none was made; every reference inside the delivered unit resolves inside it, so no unresolved reference or citation is produced. Printed slips kept literal: no printed word, symbol, hypothesis or number of the counted statement or of its proof was altered. Layout and class: the article's own class is \texttt{sn-jnl} with packages including bm, amsmath, amssymb, amsthm, mathtools, enumitem, hyperref; the wrapper is the lane's pinned \texttt{amsart} editorial wrapper at 10pt on A4, which restates the theorem-like environments the retained text needs and adds the article's own macro definitions that the retained text needs (none), with extra wrapper packages (bm, bbm, dsfont, xspace, cancel, mathtools). The delivered numbering is the unit's own. Assets: no retained span contains a figure environment, a graphics-inclusion command, a PDF-page inclusion, a table or a TikZ picture; no image file is copied into this package, the package declares no asset dependency and no image file is redistributed with it. Licence: version arXiv:2607.22478v1 of this article is distributed under the Creative Commons Attribution 4.0 International licence, as recorded in the arXiv licence metadata for this exact version and shown on the abstract page https://arxiv.org/abs/2607.22478v1. A full rights notice is delivered at licenses/arxiv-2607.22478-CC-BY-4.0.txt. Characters: every retained excerpt is pure ASCII, so the delivered engine, XeLaTeX with its default Latin Modern text font, reports no missing character.
\begin{proposition}[Differential inclusion]\label{prop:inclusion}
For $0\le p\le d-1$, we have
\begin{equation}
d\mathcal C^p\subset \mathcal C^{p+1}.
\end{equation}
More precisely, for every constant-coefficient $p$-form
$\alpha\in\Lambda^p$ and every neuron $s_i$,
\begin{equation}
d\bigl(\sigma_{k-p}(s_i)\alpha\bigr)
=
(k-p)\sigma_{k-p-1}(s_i)\,ds_i\wedge\alpha .
\end{equation}
\end{proposition}

\begin{lemma}[Induction of linear independence]\label{lem:induction}
    Let $r\in\mathbb N_0$ and suppose that
    \begin{equation}
    \{\sigma_r(s_i)\}_{i=1}^n
    \end{equation}
    is linearly independent on $\Omega$. Then for every integer $m\ge r$, the family
    \begin{equation}
    \{\sigma_m(s_i)\}_{i=1}^n
    \end{equation}
    is linearly independent on $\Omega$.
\end{lemma}

\begin{lemma}[Koszul homotopy]\label{lem:koszul}
Let $\lambda\in\Lambda^1$ be a nonzero covector. Choose a vector $v\in\mathbb R^d$ such that $\lambda(v)=1$. Then, for every $p$-form $\alpha\in\Lambda^p$,
\begin{equation}\label{eq:homo}
\iota_v(\lambda\wedge \alpha)+\lambda\wedge \iota_v\alpha=\alpha.
\end{equation}
Here $\iota_v:\Lambda^p\rightarrow\Lambda^{p-1}$ denotes contraction with the vector $v$, namely
\begin{equation}
(\iota_v\alpha)(w_1,\ldots,w_{p-1})=\alpha(v,w_1,\ldots,w_{p-1}).
\end{equation}
Furthermore, the algebraic Koszul complex
\begin{equation}
0\longrightarrow \Lambda^0
\xrightarrow{\lambda\wedge}
\Lambda^1
\xrightarrow{\lambda\wedge}
\cdots
\xrightarrow{\lambda\wedge}
\Lambda^d
\longrightarrow 0
\end{equation}
is exact.
\end{lemma}

\begin{theorem}[Exact ReLU$^k$ neural de Rham complex]\label{thm:exact}
    Assume $k\ge d$ and that the family
    \begin{equation}
    \{\sigma_{k-d}(s_i)\}_{i=1}^n
    \end{equation}
    is linearly independent on $\Omega$, then the complex
    \begin{equation}\label{eq:derham}
        0\longrightarrow\mathcal C ^0 \xrightarrow{d} \mathcal C ^1 \xrightarrow{d}
        \cdots \xrightarrow{d} \mathcal C ^d \longrightarrow 0
    \end{equation}
    is exact.
\end{theorem}

\begin{proof}
By Lemma~\ref{lem:induction}, for every $0\le p\le d$, the family
\begin{equation}
\{\sigma_{k-p}(s_i)\}_{i=1}^n
\end{equation}
is linearly independent on $\Omega$. Hence, every element of $\mathcal C^p$ has a unique representation
\begin{equation}
\nu=\sum_{i=1}^n
\sigma_{k-p}(s_i)\alpha_i,
\qquad
\alpha_i\in\Lambda^p.
\end{equation}
We first prove exactness at $\mathcal C^0$. Assume that $ d\nu=0$, Proposition~\ref{prop:inclusion} gives
\begin{equation}
0= d\nu =k\sum_{i=1}^n
\alpha_i\sigma_{k-1}(s_i)d s_i.
\end{equation}
By the linear independence of $\{\sigma_{k-1}(s_i)\}_{i=1}^n$ and the fact that the coefficient one-forms are constant, we obtain
\begin{equation}
\alpha_i d s_i=0,
\qquad i=1,\ldots,n.
\end{equation}
Since $ d s_i\neq0$, it follows that $\alpha_i=0$ for all $i$. Therefore
\begin{equation}
\nu=0.
\end{equation}
Thus
\begin{equation}
\ker( d:\mathcal C^0\to\mathcal C^1)=0,
\end{equation}
which is exactness at $\mathcal C^0$.

Now let $1\le p\le d-1$, and suppose
\begin{equation}
\nu=\sum_{i=1}^n \sigma_{k-p}(s_i)\alpha_i \in \mathcal C ^p,
\qquad d\nu=0.
\end{equation}
Again by Proposition~\ref{prop:inclusion},
\begin{equation}
0=d\nu=(k-p)\sum_{i=1}^n \sigma_{k-p-1}(s_i)\,ds_i\wedge\alpha_i.
\end{equation}
By independence of $\{\sigma_{k-p-1}(s_i)\}$,
\begin{equation}
ds_i\wedge\alpha_i=0,
\qquad i=1,\dots,n.
\end{equation}
Choose $v_i\in\mathbb  R  ^d$ such that $ds_i(v_i)=1$. Lemma~\ref{lem:koszul} gives
\begin{equation}
\alpha_i=ds_i\wedge \iota_{v_i}\alpha_i.
\end{equation}
Therefore
\begin{equation}
\nu
=\sum_{i=1}^n \sigma_{k-p}(s_i)\,ds_i\wedge \iota_{v_i}\alpha_i
=d\Bigl(\sum_{i=1}^n \frac{1}{k-p+1}\sigma_{k-p+1}(s_i)\,\iota_{v_i}\alpha_i\Bigr).
\end{equation}
Hence $\ker(d:\mathcal C ^p\to\mathcal C ^{p+1})\subset \operatorname{im}(d:\mathcal C ^{p-1}\to\mathcal C ^p)$.
The reverse inclusion is automatic because $d^2=0$.

Finally, let
\begin{equation}
\nu=\sum_{i=1}^n \sigma_{k-d}(s_i)\alpha_i\in\mathcal C ^d,
\qquad \alpha_i\in\Lambda^d.
\end{equation}
Choose $v_i$ with $ds_i(v_i)=1$. Since $\alpha_i$ has top degree, $ds_i\wedge\alpha_i=0$, so Lemma~\ref{lem:koszul} yields
\begin{equation}
\alpha_i=ds_i\wedge\iota_{v_i}\alpha_i.
\end{equation}
Hence
\begin{equation}
\nu
=d\Bigl(\sum_{i=1}^n \frac{1}{k-d+1}\sigma_{k-d+1}(s_i)\,\iota_{v_i}\alpha_i\Bigr),
\end{equation}
which proves surjectivity onto $\mathcal C ^d$.
Thus the whole complex \eqref{eq:derham} is exact.
\end{proof}

\end{document}
